\section{Software Engineering}
\label{sec:eng}

\textbf{Pure cores, thin nodes.} Every package has the same layout: \texttt{core/} holds the
algorithms, state machines and configuration dataclasses in plain Python and numpy, without
ROS imports; \texttt{utils/} holds the conversions between messages and core types; the nodes
only declare parameters, topics, services and timers. The path follower, for example, is a pure
state machine \texttt{Follower.tick(inputs)} $\to$ \texttt{Decision} (command, status, stop
events) with an explicit table of allowed state transitions; the node feeds it once per control
tick and publishes the decision, so every path through it, including zero commands on stop,
e-stop, stale inputs, pause and shutdown, is unit-tested without ROS. The original follower
node had 1429 lines; the new node has 277.

\textbf{Shared building blocks.} \pkg{dharma\_utils} provides what packages used to copy:
angles, quaternions, SE(2) and SE(3) helpers, filters (EMA, slew limiter, CUSUM), the typed
configuration, the robot profile, the braking model and swept-footprint clearances, atomic
files with backups and path guards, grid distance fields and the correlative matcher, the plugin
registry, structured logging, and \texttt{DharmaNode} (typed parameters with live validation,
TF cache with per-edge ages, exception-guarded timers, shutdown hooks).

\textbf{Checks.} A style test in every package enforces the conventions: modules of at most 300
lines (tests 400), no ROS imports in \texttt{core/}, no shadowed node members. Lint and format
run through the same pre-commit hooks locally and in CI (ruff, clang-format, prettier,
shellcheck). The CI workflow, identical in every repository, builds and tests on ROS~2 Foxy
(Python~3.8, numpy~1.17), Humble (Python~3.10, numpy~1.21) and Jazzy (Python~3.12, numpy~1.26);
Dharma dependencies are imported from \texttt{deps.repos} and checked out at the same branch when it exists, so a
change that spans repositories is tested together before it is merged.

\textbf{Traceability.} Two reviews of the original stack produced findings with priorities
(P0 safety, P1 robustness, P2 maintainability). Each one maps to a change and a regression test;
the technical manual lists the full table.

\begin{table}[t]
  \centering
  \caption{Engineering metrics before (v0.1) and after (v0.2.0)}
  \label{tab:metrics}
  \begin{tabular}{@{}lrr@{}}
    \toprule
    Metric & v0.1 & v0.2.0 \\
    \midrule
    Largest node (lines) & 1429 & 277 \\
    Copies of the node main boilerplate & 15 & 1 \\
    Copies of the latched QoS definition & 11 & 1 \\
    Robot geometry declared in (files) & $\approx 6$ & 1 \\
    Test functions (without lint) & 158 & 595 \\
    Findings with a cited test & -- & 90 of 97 \\
    CI distributions & 2 & 2 \\
    \bottomrule
  \end{tabular}
\end{table}
